\documentclass[11pt]{article}
\usepackage[margin=1in]{geometry}
\usepackage{amsmath,amssymb,amsthm,mathtools,mathrsfs}
\usepackage{microtype}
\usepackage[hidelinks]{hyperref}
\allowdisplaybreaks

\newtheorem{theorem}{Theorem}
\newtheorem{lemma}[theorem]{Lemma}
\newtheorem{proposition}[theorem]{Proposition}
\newtheorem{corollary}[theorem]{Corollary}

\newcommand{\I}{[0,1]}
\newcommand{\Fix}{\operatorname{Fix}}
\newcommand{\Turn}{\operatorname{Turn}}
\newcommand{\cl}{\operatorname{cl}}
\newcommand{\id}{\operatorname{id}}

\title{Common Fixed Points for Commuting\\Finite-Piece Piecewise-Monotone Interval Maps}
\author{}
\date{}

\begin{document}
\maketitle

\begin{abstract}
Let $P,Q:\I\to\I$ be commuting continuous maps, each monotone on finitely many
interval pieces, with constant pieces allowed. We prove that $P$ and $Q$ have a common fixed point.
After restriction and affine normalization to a minimal common invariant interval, a counterexample
has a monotone factor onto a pair of commuting expanding constant-slope PL maps. We choose such a
factor of minimum fixed-point complexity, which rules out common invariant intervals for two suitable
powers. A repelling germ then develops into a dense unit-speed fold with two independent
self-similarities. Commutation makes its joint postcritical set finite; Durand's theorem makes its finite
Markov coding eventually periodic. The resulting periodic fold has a reflection symmetry, and the
constant-slope pair lifts to commuting affine circle maps. A four-torsion calculation gives a common fixed fibre.
\end{abstract}

\section{The theorem and external inputs}

A continuous map $T:\I\to\I$ is \emph{finite-piece piecewise monotone} if it is nondecreasing or
nonincreasing on each cell of some finite interval partition. This class is closed under composition.
Throughout, products of maps denote composition from right to left: $UT=U\circ T$.

\begin{theorem}\label{thm:main}
If $P,Q:\I\to\I$ are commuting continuous finite-piece piecewise-monotone maps, then
\[
\Fix(P)\cap\Fix(Q)\ne\varnothing.
\]
\end{theorem}

Without the finite-piece hypothesis the conclusion is false: Boyce and Huneke independently constructed
commuting continuous interval maps with no common fixed point \cite{Boyce,Huneke}. The reduction begins
with two published inputs. Bogatyi--Frolkina prove that commuting continuous selfmaps
$F,G$ of an interval with no common fixed point satisfy
\begin{equation}\label{eq:BF}
h_{\rm top}(FG)\ge 2\log2>0
\end{equation}
(with natural logarithms) \cite[Theorem~1]{BogatyiFrolkina}. Second, Alsed\`a--Misiurewicz's constant-slope theorem gives the following eigenmetric consequence
of their Theorem~A and Lemmas~1.1 and~3.1 \cite[Theorem~A and Lemmas~1.1, 3.1]{AlsedaMisiurewicz}: if a
finite-piece piecewise-monotone map $H$ has entropy $\log\lambda>0$, there is a continuous
nondecreasing surjection $\psi:\I\to\I$ such that on every monotonicity interval of $H$,
\begin{equation}\label{eq:AM}
|\psi(Hy)-\psi(Hx)|=\lambda|\psi(y)-\psi(x)|.
\end{equation}
Indeed, Theorem~A supplies an $H$-compatible increasing semiconjugacy to a map of constant slope
$\lambda$. Lemma~1.1 says that the constant-slope factor is monotone on the image of each
monotonicity interval of $H$; on a nondegenerate image its slope has absolute value $\lambda$, and
on a degenerate image both sides above vanish. This also covers non-strict monotonicity.

\section{Reduction to a minimal constant-slope counterexample}

A continuous map is \emph{finite-piece PL} if it is affine on every cell of some finite interval partition, and
\emph{flat-free} if it is constant on no nondegenerate interval. A flat-free finite-piece PL map has
\emph{constant absolute slope} $c>0$ if every affine branch has slope $+c$ or $-c$. A
\emph{common monotone factor} of $P,Q$ is given by a continuous nondecreasing surjection
$\pi:\I\to\I$ and maps $\bar P,\bar Q$ satisfying
$\pi P=\bar P\pi$ and $\pi Q=\bar Q\pi$. We call a counterexample \emph{minimal} if it has no
proper nonempty compact interval which is forward invariant under both generators.

For a finite-piece piecewise-monotone map $T$ and $S\in C(\I,\mathbb R)$ with $S(0)=0$, define the
\emph{branch length} $\mathcal L_TS$ by adding the $S$-lengths of the images of successive
monotonicity pieces. Explicitly, if $0=c_0<\cdots<c_r=1$ is a monotonicity partition and
$x\in[c_{j-1},c_j]$,
\[
(\mathcal L_TS)(x)=\sum_{i<j}[S;T(c_{i-1}),T(c_i)]+[S;T(c_{j-1}),T(x)],
\]
where $[S;u,v]=S(\max\{u,v\})-S(\min\{u,v\})$.

\begin{lemma}[Common constant-slope factor]\label{lem:linearize}
Every counterexample to Theorem~\ref{thm:main} has a nondegenerate compact common invariant
interval $J$ which is minimal among such intervals. After affinely conjugating the restrictions
$P|_J,Q|_J$ to maps of $\I$, the resulting commuting onto minimal-core counterexample admits a
common monotone factor $(\bar P,\bar Q)$. The factor maps are commuting onto finite-piece PL maps of
constant absolute slopes $C,D>1$; they have no common fixed point and no proper nonempty compact
common invariant interval.
\end{lemma}

\begin{proof}
We use the following elementary properties of branch length. The displayed formulas agree at the endpoints
of adjacent partition cells, so $\mathcal L_TS$ is continuous and vanishes at $0$. Refining a
monotonicity partition does not change $\mathcal L_T$, because splitting a monotonicity interval at
$d$ splits one $S$-length into two. The operator $\mathcal L_T$ is linear, and it sends
nondecreasing functions to nondecreasing functions:
if $x<y$ lie in one monotonicity interval of $T$, then
\begin{equation}\label{eq:length-increment}
(\mathcal L_TS)(y)-(\mathcal L_TS)(x)=[S;T(x),T(y)],
\end{equation}
which is nonnegative when $S$ is nondecreasing. Finally, after refining a partition of $T$ so that every
image cell lies in one monotonicity cell of $U$, \eqref{eq:length-increment} gives
\begin{equation}\label{eq:length-composition}
\mathcal L_{UT}=\mathcal L_T\mathcal L_U.
\end{equation}

We first pass to a minimal common invariant interval. Let $\mathscr J$ be the set of pairs $(a,b)$
with $0\le a\le b\le1$ for which both $P([a,b])$ and $Q([a,b])$ are contained in $[a,b]$.
It is nonempty because $(0,1)\in\mathscr J$. It is closed: if
$(a_k,b_k)\to(a,b)$ with $(a_k,b_k)\in\mathscr J$ and $x\in[a,b]$, choose
$x_k\in[a_k,b_k]$ with $x_k\to x$. Continuity gives
$P(x)=\lim P(x_k)\in[a,b]$, and likewise for $Q$. Thus $\mathscr J$ is a closed subset of the
compact triangle $0\le a\le b\le1$; choose $(a,b)\in\mathscr J$ of least length and put
$J=[a,b]$. The interval $J$ is nondegenerate, since an invariant singleton would be a common fixed
point, and least length makes it inclusion-minimal.

The image $P(J)$ is a nonempty compact interval contained in $J$, and
\[
P(P(J))\subseteq P(J),\qquad Q(P(J))=P(Q(J))\subseteq P(J).
\]
Thus $P(J)$ is another compact common invariant interval, so minimality gives $P(J)=J$; similarly
$Q(J)=J$. Let $\eta:\I\to J$ be the increasing affine homeomorphism and replace
$P|_J,Q|_J$ by $\eta^{-1}P\eta,\eta^{-1}Q\eta$. These normalized maps are commuting and onto,
have no common fixed point, and have no proper nonempty compact common invariant interval. We retain
the symbols $P,Q$ for them. The map $\eta$ conjugates this normalized restriction to the subsystem
on $J$; it does not make the restriction a factor of the original global system. The monotone
coordinate constructed below is a factor of the normalized pair. Put $H=PQ$. By \eqref{eq:BF},
$h_{\rm top}(H)>0$; write $\lambda=e^{h_{\rm top}(H)}$.

Let
\[
\mathscr E=\{S\in C(\I,\mathbb R):S(0)=0,\ \mathcal L_HS=\lambda S\}.
\]
This space is finite-dimensional. Indeed, choose $N$ with $\lambda^N>2$ and a monotonicity partition
$0=c_0<\cdots<c_r=1$ for $H^N$, and set
\[
V_N=\{H^N(c_0),\ldots,H^N(c_r)\}.
\]
We claim that evaluation on the finite set $V_N$ is injective on $\mathscr E$. Suppose that
$S\in\mathscr E$ vanishes on $V_N$, and put $M=\|S\|_\infty$. Since
$\mathcal L_{H^N}=\mathcal L_H^N$, for $x\in[c_{j-1},c_j]$ the definition of branch length gives
\[
\lambda^NS(x)
 =\sum_{i<j}[S;H^N(c_{i-1}),H^N(c_i)]
  +[S;H^N(c_{j-1}),H^N(x)].
\]
Every completed-lap term vanishes because both endpoints lie in $V_N$, while the last bracket has one
endpoint in $V_N$. Thus $|\lambda^NS(x)|\le M$.
Taking the supremum gives $M\le M/\lambda^N$, hence $M=0$. Evaluation on $V_N$ is injective, so
$\dim\mathscr E\le\#V_N<\infty$.

Since $\psi$ is nondecreasing, \eqref{eq:length-increment} and \eqref{eq:AM} give
$\mathcal L_H\psi=\lambda\psi$; thus $\psi\in\mathscr E$. Let $A=\mathcal L_P$ and
$B=\mathcal L_Q$. Since $P$ and $Q$ commute with $H$, \eqref{eq:length-composition} gives
$\mathcal L_HA=A\mathcal L_H$ and $\mathcal L_HB=B\mathcal L_H$; hence $A$ and $B$ preserve
$\mathscr E$. The same composition rule and $PQ=QP=H$ give
\begin{equation}\label{eq:AB}
AB=BA=\lambda\id\qquad\text{on }\mathscr E.
\end{equation}
Consider
\[
K=\{S\in\mathscr E:S\text{ is nondecreasing},\ S(1)=1\}.
\]
Because $\psi$ is a nondecreasing surjection, $\psi(0)=0$ and $\psi(1)=1$, so $K$ contains $\psi$. It is closed and convex, and every $S\in K$ satisfies $0\le S\le1$; since
$\mathscr E$ is finite-dimensional, $K$ is compact. If $0=d_0<\cdots<d_s=1$ is a monotonicity partition of the onto map $P$, then
for $S\in K$,
\[
(AS)(1)=\sum_i[S;P(d_{i-1}),P(d_i)]\ge1,
\]
because the intervals $P([d_{i-1},d_i])$ cover $\I$, so their $S$-images cover $[S(0),S(1)]=[0,1]$ and the sum of their lengths is at least $1$. Thus $S\mapsto AS/(AS)(1)$ is a continuous
selfmap of $K$. Brouwer gives $\theta\in K$ and $C>0$ with $A\theta=C\theta$; by
\eqref{eq:AB}, $B\theta=D\theta$ with $D=\lambda/C$. Indeed, \eqref{eq:AB} makes $A$ injective,
and $A(B\theta-D\theta)=0$.

If $\theta(x)=\theta(y)$ for $x<y$, subdivide $[x,y]$ at the monotonicity endpoints of $P$.
Equation \eqref{eq:length-increment} gives zero $\theta$-length to every image subinterval, hence
$\theta(Px)=\theta(Py)$; similarly for $Q$. Thus both maps descend through $\theta$. On the image
of any monotonicity interval on which $\theta$ is nonconstant, \eqref{eq:length-increment} gives
\[
|\bar P(t)-\bar P(s)|=C|t-s|,
\]
and similarly for $\bar Q$. Hence the descendants are finite-piece PL of constant absolute slopes
$C,D$. Indeed, the nondegenerate $\theta$-images of the finitely many monotonicity cells cover $\I$;
their finitely many endpoints determine a common refinement, and every refined cell lies in one such
image interval, where the displayed affine formula holds. Since $C,D>0$, neither descendant has a flat
interval. They are onto because the original restrictions $P,Q$ are onto and $\theta$ is surjective.

A nonconstant monotone factor of a minimal counterexample is again a minimal counterexample. More
precisely, if $\pi$ is such a factor and $T$ is any word in $P,Q$, then
\begin{equation}\label{eq:fixed-descent}
\Fix(\bar T)=\pi(\Fix(T)).
\end{equation}
The inclusion from right to left is immediate. Conversely, if $\bar T(y)=y$, the interval fibre
$\pi^{-1}(y)$ is $T$-invariant, so the intermediate value theorem gives a fixed point of $T$ in that
fibre. If a point $y$ were fixed by both descended maps, then $\pi^{-1}(y)$ would be a nonempty compact
interval invariant under both original maps; it is proper because $\pi$ is nonconstant, contradicting
minimality. Likewise, the inverse image of a proper nonempty compact common invariant interval downstairs
would be a proper nonempty compact common invariant interval upstairs. Thus the factor has neither a
common fixed point nor a proper common invariant interval, and is again a minimal counterexample.

Finally $C,D>1$. For an onto flat-free constant-slope map the total variation equals its absolute
slope, so the slope is at least $1$. Equality makes the map monotone, hence either $x\mapsto x$ or
$x\mapsto1-x$; in the first case any fixed point of the other map is common, and in the second
commutation forces the other map to fix $1/2$. Both alternatives contradict counterexamplehood.
\end{proof}

We shall also use Fugate--McLean's theorem that the fixed-point set of a periodic homeomorphism of a
tree-like continuum is a nonempty continuum \cite[Theorem~1.5, p.~613]{FugateMcLean}.

\begin{lemma}[Common positive iterates]\label{lem:common-iterate}
If commuting continuous interval maps $F,G$ satisfy $F^m=G^n$ for some $m,n\ge1$, then $F,G$ have a
common fixed point.
\end{lemma}

\begin{proof}
Put $T=F^m=G^n$ and let
\[
J_\infty=\bigcap_{r\ge0}T^r(\I).
\]
The intervals $T^r(\I)$ are nested nonempty compact intervals, so $J_\infty$ is a nonempty compact
interval. Also $T(J_\infty)=J_\infty$: the forward inclusion is immediate, while for
$y\in J_\infty$ the nested compact sets $T^{-1}(y)\cap T^r(\I)$ are nonempty and have a common
point in $J_\infty$. Commutation gives
\[
F(J_\infty)\subseteq J_\infty,\qquad G(J_\infty)\subseteq J_\infty.
\]
If $J_\infty$ is a singleton, both $F$ and $G$ preserve that singleton, so its point is a common
fixed point.

Assume therefore that $J_\infty$ is nondegenerate and form the natural extension
\[
X_T=\varprojlim(J_\infty,T)
 =\{(x_0,x_1,\ldots):T(x_{j+1})=x_j\}.
\]
For each $N$, the subset of $J_\infty^{\mathbb N_0}$ obtained by imposing only the first $N$ bonding
equations is nonempty, compact, and connected: projection to the coordinates
$x_N,x_{N+1},\ldots$ identifies it with a product of copies of $J_\infty$. These subsets are nested,
and their intersection is $X_T$. Thus $X_T$ is a nonempty compact connected metric space. Because
$T|_{J_\infty}$ is onto, every coordinate projection $\pi_N:X_T\to J_\infty$ is onto. Give $X_T$
the metric
\[
d_X(x,y)=\sum_{i=0}^{\infty}2^{-i-1}|x_i-y_i|.
\]
If $\pi_N(x)=\pi_N(y)$, the bonding equations give $x_i=y_i$ for $0\leq i\leq N$, and hence
\[
d_X(x,y)\leq\sum_{i>N}2^{-i-1}=2^{-N-1}.
\]
The projections $\pi_N$ are therefore maps onto an interval with arbitrarily small fibres. Hence
$X_T$ is arc-like, in particular chainable and tree-like. The shift
\[
\Sigma(x_0,x_1,x_2,\ldots)=(T(x_0),x_0,x_1,\ldots)
\]
is a homeomorphism.

Coordinatewise application of $F$ and $G$ defines commuting lifts $\widehat F,\widehat G:X_T\to X_T$.
Because $F^m=G^n=T$,
\[
\widehat F^{\,m}=\Sigma=\widehat G^{\,n},
\]
so both lifts are homeomorphisms: if a positive power of a selfmap is bijective, then the selfmap itself
is injective and surjective, and a continuous bijection of a compact Hausdorff space is a homeomorphism.
Write $m=dm_0$, $n=dn_0$ with
$\gcd(m_0,n_0)=1$, choose $r,s\in\mathbb Z$ with $rn_0+sm_0=1$, and set
\[
U=\widehat F^{\,r}\widehat G^{\,s},\qquad
V=\widehat F^{\,m_0}\widehat G^{-n_0}.
\]
Then
\[
V^d=\id,\qquad
\widehat F=U^{n_0}V^s,\qquad
\widehat G=U^{m_0}V^{-r}.
\]
Put $\mathcal D=\Fix(V)$. The quoted Fugate--McLean
theorem says that $\mathcal D$ is a nonempty continuum. Since $U$ commutes with $V$, it preserves
$\mathcal D$.
We prove directly that $U|_{\mathcal D}$ has a fixed point. Given $\varepsilon>0$, choose $N$ with
$2^{-N-1}<\varepsilon$. Connectedness gives
$\pi_N(\mathcal D)=[\alpha,\beta]$. If $\alpha=\beta$, every $z\in\mathcal D$ satisfies
$\pi_N(Uz)=\pi_N(z)$. If $\alpha<\beta$, choose $z_\alpha,z_\beta\in\mathcal D$ over the two
endpoints. The continuous function
\[
z\longmapsto \pi_N(Uz)-\pi_N(z)
\]
is nonnegative at $z_\alpha$ and nonpositive at $z_\beta$; its image of the connected set
$\mathcal D$ therefore contains zero. In either case there is $z_\varepsilon\in\mathcal D$ with
$\pi_N(Uz_\varepsilon)=\pi_N(z_\varepsilon)$. The bonding equations then give equality in every
coordinate $0,\ldots,N$, and hence
\[
d_X(Uz_\varepsilon,z_\varepsilon)\leq
\sum_{i>N}2^{-i-1}=2^{-N-1}<\varepsilon.
\]
Take $\varepsilon_j\downarrow0$ and pass to a convergent subsequence of $z_{\varepsilon_j}$ in the
compact set $\mathcal D$. Continuity of $U$ shows that its limit $z$ satisfies $U(z)=z$. Since also
$V(z)=z$, the displayed Bezout factorizations show that
$\widehat F(z)=z=\widehat G(z)$. The zeroth coordinate of $z$ is a common fixed point of $F$ and
$G$.
\end{proof}

\begin{lemma}[Extremal irreducibility]\label{lem:irreducible}
Let $P_0,Q_0$ be a minimal constant-slope counterexample supplied by Lemma~\ref{lem:linearize}. Put
\[
\nu_0=\#\Fix(P_0Q_0),\qquad L=(2\nu_0)!.
\]
The integer $\nu_0$ is positive by the interval fixed-point theorem and finite because $P_0Q_0$ is finite-piece PL with branch-slope magnitude $>1$; in particular $L\geq1$.
Then $P_0,Q_0$ have a constant-slope minimal-counterexample factor $P,Q$ such that, with
$A=P^L$ and $B=Q^L$,
\begin{equation}\label{eq:irreducible}
A(J)=J=B(J),\quad J\subseteq\I\text{ a nondegenerate compact interval}
\quad\Longrightarrow\quad J=\I.
\end{equation}
Moreover $\#\Fix(PQ)\le\nu_0$.
\end{lemma}

\begin{proof}
The class of constant-slope minimal-counterexample factors of $(P_0,Q_0)$ is nonempty, since it contains
$(P_0,Q_0)$. For such a factor set
\[
\mu(P,Q)=\#\Fix(P^{2L}).
\]
This is a positive finite integer: existence of a fixed point follows from the interval fixed-point theorem,
and finiteness follows because $P^{2L}-\id$ is piecewise affine with nonzero slope on every affine cell
(the branch slopes of $P^{2L}$ have magnitude $C^{2L}>1$). By the well-ordering of the positive integers, choose such a factor $(P,Q)$ with $\mu(P,Q)$ minimal.
Equation \eqref{eq:fixed-descent} gives $\#\Fix(PQ)\le\nu_0$.

Suppose a proper nondegenerate $J$ satisfies \eqref{eq:irreducible}. For
$D_{i,j}=P^iQ^j(J)$, $(i,j)\in(\mathbb Z/L\mathbb Z)^2$, form the intersection graph. Translation in
either index is an automorphism: forward translation preserves intersections, and the converse follows
by applying $P^{L-1}$ or $Q^{L-1}$ and using $P^L(D_{i,j})=D_{i,j}=Q^L(D_{i,j})$.

If the graph is connected, it is complete. Indeed, in any finite family of intervals, one with least
right endpoint is simplicial; vertex-transitivity makes every vertex simplicial, whereas a shortest path
of length at least two has an interior vertex with two nonadjacent neighbours. Thus all $D_{i,j}$ meet
pairwise, and Helly's property gives a nonempty common intersection. That intersection is forward
invariant under $P,Q$, so minimality makes it $\I$, contradicting $J\ne\I$.

Hence the graph is disconnected. The union of the intervals in each graph component is a nondegenerate
compact interval, and distinct component-unions are disjoint. Write them in their order on the interval as
\[
K_r=[\alpha_r,\beta_r]\qquad(1\leq r\leq N),
\quad \beta_r<\alpha_{r+1}.
\]
There are at least two of them. Translation of the phase labels is a graph automorphism and satisfies
$P(D_{i,j})=D_{i+1,j}$ and $Q(D_{i,j})=D_{i,j+1}$. Consequently $P$ and $Q$ permute the intervals
$K_r$.

Put
\[
\kappa=1-\sum_{r=1}^N(\beta_r-\alpha_r)>0.
\]
The strict inequality follows from the nonempty gaps between the distinct component intervals. If
$|K_r\cap[0,x]|$ denotes the length of that intersection, define
\[
\pi(x)=\frac{1}{\kappa}
 \left(x-\sum_{r=1}^N|K_r\cap[0,x]|\right).
\]
Then $\pi:\I\to\I$ is a continuous nondecreasing PL surjection, constant on each $K_r$ and strictly
increasing on every complementary gap. Its only non-singleton fibres are the intervals $K_r$. Since each
generator sends every $K_r$ onto one of the $K_s$, both generators respect the fibres of $\pi$. There are
therefore unique continuous maps $\bar P,\bar Q:\I\to\I$ such that
\[
\pi P=\bar P\pi,\qquad \pi Q=\bar Q\pi.
\]
Surjectivity of $P,Q$ and $\pi$ makes $\bar P,\bar Q$ onto, and surjectivity of $\pi$ transfers the
commutation relation to them.

It remains to check finite-piece PL regularity. On each complementary gap, $\pi$ has an affine inverse.
Refine that gap at the affine breakpoints of $P$ and $Q$ and at the preimages of all endpoints
$\alpha_r,\beta_r$; there are only finitely many such points because the generators are flat-free
finite-piece PL maps. On a resulting cell $G$, a generator image lies either in one $K_r$, in which case
the descended map is constant, or in a single complementary gap, in which case
\[
\bar P|_{\pi(G)}=\pi\circ P\circ(\pi|_G)^{-1}
\quad\text{or}\quad
\bar Q|_{\pi(G)}=\pi\circ Q\circ(\pi|_G)^{-1}
\]
is affine. Thus both descendants are finite-piece PL. The factor and fixed-point arguments from
Lemma~\ref{lem:linearize}, including \eqref{eq:fixed-descent}, show that the descended pair is again a
minimal counterexample.

This collapse strictly lowers $\#\Fix(P^{2L})$. The map $A|_J$ is onto, and every continuous onto
selfmap $f$ of a nondegenerate interval has at least two fixed points of $f^2$. To see this, write the
interval as $[u,v]$. If $f$ already has two fixed points there is nothing to prove. Otherwise let $z$ be
its unique fixed point. It lies in $(u,v)$: if, for example, $z=u$, choose $x$ with $f(x)=v$;
if $x=v$ then $v$ is another fixed point, while if $x<v$ the signs of $f(x)-x>0$ and
$f(v)-v\le0$ give another zero by the intermediate value theorem. The case $z=v$ is symmetric.
Thus $z\in(u,v)$, and uniqueness of the zero of $f(x)-x$, together with
$f(u)-u\ge0$ and $f(v)-v\le0$, gives
$f(x)>x$ for $x<z$ and $f(x)<x$ for $x>z$. Surjectivity gives $a<z<b$ with $f(a)=v$ and $f(b)=u$: a preimage of $v$ cannot lie to the
right of $z$, where $f(x)<x\le v$, and a preimage of $u$ cannot lie to the left of $z$, where
$f(x)>x\ge u$. Since $f(a)=v\ge b$ while $f(z)=z<b$, the intermediate value theorem gives
$c\in[a,z)$ with $f(c)=b$. Consequently
\[
f^2(c)=u\le c,\qquad f^2(u)\ge u.
\]
If equality holds at either endpoint, that endpoint is fixed by $f^2$; otherwise the intermediate value
theorem gives a zero of $f^2(x)-x$ in $(u,c)$. In either case $f^2$ has a fixed point strictly to the
left of $z$, while $z$ itself is fixed by $f^2$. Applied to $f=A|_J$, this gives two distinct
points of $J\cap\Fix(A^2)=J\cap\Fix(P^{2L})$, and the collapse identifies them. By
\eqref{eq:fixed-descent}, the descended $P^{2L}$ therefore has strictly fewer fixed points.

The quotient pair need not have constant slopes. Apply Lemma~\ref{lem:linearize} to it. Because the
quotient pair is already minimal on its common invariant core, the preliminary minimal-core restriction
in that lemma is the whole interval. The resulting constant-slope system is therefore a monotone factor
of the quotient pair. Composing this coordinate with the collapse coordinate and with the original
factor coordinate shows that the resulting pair still belongs to the class of constant-slope
minimal-counterexample factors of $(P_0,Q_0)$. Both slopes are greater than one, and
\eqref{eq:fixed-descent} for the composite factor coordinate shows that the number of fixed points of
$P^{2L}$ cannot increase. This contradicts the defining minimality of $\mu$ and proves
\eqref{eq:irreducible}.
\end{proof}

\section{The two-scale fold}

Fix the extremal pair $P,Q$. Write its slopes as $C,D>1$ and set
\begin{equation}\label{eq:scales}
A=P^L,\quad B=Q^L,\quad R=AB,\qquad a=C^L,\quad b=D^L,\quad s=ab.
\end{equation}

For a flat-free finite-piece PL map $T$, an interior point $x$ is a \emph{turn} if, on sufficiently
small left- and right-hand intervals at $x$, the map $T$ is affine with nonzero slopes of opposite
signs. This definition is intrinsic: inserting a redundant vertex into an affine branch does not create
a turn. We write
\[
E(T)=\{0,1\}\cup\Turn(T),\qquad \mathcal V(T)=T(E(T)).
\]
Two numbers $a,b>1$ are \emph{multiplicatively independent} if
$a^m\ne b^n$ for every pair of positive integers $m,n$.

\begin{lemma}[Finite joint postcritical set]\label{lem:finite-postcritical}
Assume that the slope magnitudes $a,b>1$ of the commuting onto flat-free finite-piece PL maps $A,B$
are multiplicatively independent. Then
\begin{equation}\label{eq:postcritical}
\mathcal C=\bigcup_{m,n\ge0}A^mB^n(E(A)\cup E(B))
\end{equation}
is finite.
\end{lemma}

\begin{proof}
We begin with the vertex-value identity
\begin{equation}\label{eq:vertices}
\mathcal V(UT)=\mathcal V(U)\cup U(\mathcal V(T))
\end{equation}
for onto flat-free finite-piece PL maps $U,T$.

For one inclusion, let $x\in E(UT)$. If $x\in E(T)$, then
$UT(x)\in U(\mathcal V(T))$. Otherwise $x$ is an interior point at which $T$ is locally one-to-one;
if $UT$ turns at $x$, then $U$ must turn at $T(x)$, so $UT(x)\in\mathcal V(U)$. Thus the left side
of \eqref{eq:vertices} is contained in the right side.

Conversely, an interior turn of $T$ remains a turn after composition with $U$: the two adjacent branches
of $T$ approach $T(x)$ from the same side, and flat-freeness makes $U$ nonconstant on that one-sided
neighbourhood. Domain endpoints of $T$ are already endpoints of $UT$. Hence
$U(\mathcal V(T))\subseteq\mathcal V(UT)$. Now let $y\in E(U)$. If $y\in(0,1)$, onto-ness of $T$
provides a locally one-to-one preimage $x$ of $y$: indeed, take a boundary point between the nonempty
open sets $\{T<y\}$ and $\{T>y\}$; a flat-free PL map crosses the level there and so does not turn.
If $y$ is a turn of $U$, then $UT$ turns at this $x$, giving $U(y)\in\mathcal V(UT)$. If
$y\in\{0,1\}$, any preimage of $y$ under $T$ is either a domain endpoint or an interior extremum, hence
belongs to $E(T)$, and the same conclusion follows. This proves \eqref{eq:vertices}.

Set
\[
\mathcal T_A(m)=\bigcup_{0\le i<m}A^i\mathcal V(A),\qquad
\mathcal T_B(n)=\bigcup_{0\le j<n}B^j\mathcal V(B).
\]
Iterating \eqref{eq:vertices} and comparing the two expressions for $A^mB^n=B^nA^m$ gives
\begin{equation}\label{eq:rectangle}
\mathcal T_A(m)\cup A^m\mathcal T_B(n)
 =\mathcal T_B(n)\cup B^n\mathcal T_A(m).
\end{equation}
Taking $n=1$ yields
\begin{equation}\label{eq:difference}
B(\mathcal T_A(m))\setminus\mathcal T_A(m)\subseteq A^m\mathcal V(B).
\end{equation}

We claim that $\mathcal T_A=\bigcup_{i\ge0}A^i\mathcal V(A)$ is finite. If it were infinite, at least one
of the finitely many seeds in $\mathcal V(A)$ would have an infinite forward $A$-orbit. Among points on
infinite such orbits, call two points tail-equivalent when some forward iterate of one equals some forward
iterate of the other. The finitely many seeds determine only finitely many infinite tail classes.

If $w$ lies on an infinite tail, then $B(w)$ also has infinite $A$-orbit. The map $B$ is finite-to-one,
because each of its finitely many nondegenerate affine branches is injective. Since
$A^jB(w)=B(A^j w)$, a finite $A$-orbit of $B(w)$ would put the infinite set
$\{A^j w:j\ge0\}$ in the preimage under $B$ of a finite set. We claim that
$B(w)\in\mathcal T_A$. Choose $m_0$ with $w\in\mathcal T_A(m_0)$. If not, then
\eqref{eq:difference} gives
$B(w)\in A^m\mathcal V(B)$ for every $m\ge m_0$. Since $\mathcal V(B)$ is finite, there are
$m_1<m_2$ and $v\in\mathcal V(B)$ with
$B(w)=A^{m_1}v=A^{m_2}v$, making $B(w)$ $A$-periodic, a contradiction. Commutation also preserves
tail-equivalence, so $B$ induces a selfmap of the finite set of infinite tail classes. Choose a periodic
class, say of period $r\ge1$, and $w$ in that class. Then for some $p_0,q_0\ge0$,
\begin{equation}\label{eq:tail-relation}
A^{p_0}B^r(w)=A^{q_0}(w).
\end{equation}
By commutation, the two finite-piece affine maps $A^{p_0}B^r$ and $A^{q_0}$ agree at every point of the
infinite forward $A$-orbit of $w$. A common finite affine partition has a cell containing infinitely many
of those distinct orbit points, so the two affine restrictions on that cell coincide. Their slope magnitudes
therefore satisfy
\[
a^{p_0}b^r=a^{q_0}.
\]
If $q_0\le p_0$, the left side after division by $a^{q_0}$ is $>1$; if $q_0>p_0$, then
$b^r=a^{q_0-p_0}$, contradicting multiplicative independence. Hence $\mathcal T_A$ is finite. By
symmetry, so is $\mathcal T_B$.

Choose $m_A$ with $\mathcal T_A(m)=\mathcal T_A$ for $m\ge m_A$. If
$y\in B(\mathcal T_A)\setminus\mathcal T_A$, then \eqref{eq:difference} gives
$y\in A^m\mathcal V(B)$ for every $m\ge m_A$; the same repeated-seed argument makes $y$ periodic
under $A$. Since there are only finitely many such $y$, some $r_A\ge1$ fixes all of them under
$A^{r_A}$. The set $\Fix(A^{r_A})$ is finite because $A^{r_A}$ is finite-piece PL with slope magnitude
$a^{r_A}>1$, and
\[
\mathcal T_A\cup\Fix(A^{r_A})
\]
is forward invariant under both $A$ and $B$: $A$ preserves each part, $B(\mathcal T_A)$ lies in their
union, and $B$ preserves $\Fix(A^{r_A})$ by commutation. Thus the semigroup orbit of $\mathcal T_A$
is finite. Symmetrically, the semigroup orbit of $\mathcal T_B$ is finite.

Finally,
\begin{equation}\label{eq:turn-transport}
B(E(A))\subseteq E(A)\cup\mathcal T_B,
\qquad
A(E(B))\subseteq E(B)\cup\mathcal T_A.
\end{equation}
For example, let $x$ be an interior turn of $A$. If $x$ is also a turn of $B$, then
$B(x)\in\mathcal V(B)\subseteq\mathcal T_B$. Otherwise $B$ is locally one-to-one at $x$. The turn of
the inner map $A$ survives in $BA=AB$; viewed as $AB$, that turn must therefore come from a turn of
the outer map $A$ at $B(x)$, so $B(x)\in E(A)$. If $x$ is an endpoint, then
$B(x)\in\mathcal V(B)$. This proves the first inclusion, and the second is symmetric. Since
$A(E(A))\subseteq\mathcal T_A$ and $B(E(B))\subseteq\mathcal T_B$, the semigroup orbit in
\eqref{eq:postcritical} is contained in the union of the finite sets $E(A),E(B)$ and the finite semigroup
orbits of $\mathcal T_A,\mathcal T_B$. Hence $\mathcal C$ is finite.
\end{proof}

A map $\Phi:[0,\infty)\to\I$ is \emph{locally finite-piece affine of unit speed} if every compact
parameter interval admits a finite subdivision on whose cells $\Phi$ is affine with slope $+1$ or
$-1$. It is \emph{eventually periodic} if there are $T_0\geq0$ and $\tau>0$ such that
$\Phi(t+\tau)=\Phi(t)$ for every $t\geq T_0$.

\begin{lemma}[Two-scale rigidity]\label{lem:two-scale}
After possibly conjugating $P,Q$ simultaneously by the reflection $x\mapsto1-x$ and redefining
$A,B,R$ accordingly, there are $p\in\I$ and a continuous locally finite-piece affine map
$\Psi:[0,\infty)\to\I$ such that every affine piece has slope $\pm1$, $\Psi(t)=p+t$ near $0$, and
\begin{equation}\label{eq:equivariance}
\Psi(at)=A\Psi(t),\qquad \Psi(bt)=B\Psi(t),\qquad \Psi(st)=R\Psi(t).
\end{equation}
Its image is dense, $a,b$ are multiplicatively independent, and $\Psi$ is eventually periodic.
\end{lemma}

\begin{proof}
Put $H=PQ$ and $X=\Fix(H)$. The set $X$ is nonempty by the interval fixed-point theorem, and it is
finite because $H$ is finite-piece PL of slope magnitude $CD>1$. Since $P,Q$ commute with $H$, they
preserve $X$, and on $X$ the relation $QP=PQ=H$ gives mutually inverse permutations
$P|_X$ and $Q|_X$. If $X$ were a singleton, both permutations would fix its unique point, giving a common
fixed point of $P,Q$. Hence $2\le\#X$. Finally \eqref{eq:fixed-descent} along the factor chain gives
$\#X\le\nu_0$.

Choose consecutive $p<c$ in $X$. Since $H(x)-x$ has no zero in $(p,c)$, it has a constant strict
sign there. If that sign is negative, let $\iota(x)=1-x$ and replace both generators simultaneously by
\[
\widetilde P=\iota P\iota,\qquad \widetilde Q=\iota Q\iota.
\]
For every word $W$ in the two generators one has $\widetilde W=\iota W\iota$. Thus this conjugacy
preserves continuity, onto-ness, finite-piece PL structure, constant slope magnitudes, commutation, the
absence of a common fixed point, all fixed-set cardinalities, and minimality of the common invariant
core. It also carries nondegenerate common invariant intervals for the returned powers bijectively to
such intervals, so \eqref{eq:irreducible} is preserved. Finally,
\[
(\iota H\iota)(x)-x=-\bigl(H(1-x)-(1-x)\bigr),
\]
and the reflected gap is $(1-c,1-p)$. After this simultaneous conjugacy and a relabelling, we may
therefore assume $H(x)>x$ on $(p,c)$. Since $H(p)=p$ and the right-hand affine germ of $H$ has
slope magnitude $CD$, the sign condition forces
$H(p+t)=p+CDt$ for all sufficiently small $t>0$.

Represent a directed side by $(x,\sigma)$, where $x\in X$, $\sigma\in\{-1,+1\}$, and
$x+\sigma t\in\I$ for all sufficiently small $t>0$. We say that $T$ carries $(x,\sigma)$ to
$(T(x),\tau)$ if there are $c,\delta>0$ such that
\[
T(x+\sigma t)=T(x)+\tau ct\qquad(0<t<\delta).
\]
For a flat-free PL map the germ on each source side is affine with nonzero slope, so it carries that
side to a unique target side. Shrinking source intervals shows that side transport respects composition.

Call a side repelling when $H$ fixes it and acts on it by positive dilation of factor $CD$. The identity
$PH=HP$ transports such a germ to a repelling side based at $P(x)$; the same is true for $Q$. Hence
$P$ and $Q$ induce well-defined maps $P_*,Q_*$ on the finite set of repelling sides. Composition
compatibility and $PQ=QP=H$ give
\[
Q_*P_*=H_*=\id=P_*Q_*.
\]
Thus the two side actions are inverse permutations. If there are $N$ repelling sides, then
$1\leq N\leq2\#X\leq2\nu_0$. Every cycle length is at most $N$, so the order of either permutation
divides $N!$, which divides $(2\nu_0)!=L$. The returned maps $A=P^L$ and $B=Q^L$ therefore fix the
chosen right side. Their slope magnitudes are $a,b$, and preservation of its orientation gives the
positive germs
\begin{equation}\label{eq:germs}
A(p+t)=p+at,\qquad B(p+t)=p+bt
\end{equation}
for all sufficiently small $t\ge0$.

Choose $\varepsilon>0$ so that \eqref{eq:germs} holds on $[p,p+s\varepsilon]$, and set
$I_n=R^n([p,p+\varepsilon])$. Since $s>a,b>1$,
$I_0\subseteq R(I_0)$ and $A(I_0),B(I_0)\subseteq R(I_0)$. Thus the $I_n$ are nested. Commutation and $A(I_0),B(I_0)\subseteq I_1$ give
$A(I_n),B(I_n)\subseteq I_{n+1}$ for every $n$, while
$R(I_n)=I_{n+1}$. Let $U=\bigcup_{n\geq0}I_n$ and $K_0=\cl U$. The sets $I_n$ are nested compact
intervals and $I_0$ has positive length, so $K_0$ is a nondegenerate compact interval. The inclusions
above give $A(U),B(U)\subseteq U$, and continuity gives
$A(K_0),B(K_0)\subseteq K_0$. Moreover,
\[
R(U)=\bigcup_{n\geq0}I_{n+1}=U,
\]
where the last equality uses $I_0\subseteq I_1$. Hence continuity gives
$R(K_0)\subseteq\cl R(U)=K_0$. Conversely, $R(K_0)$ is compact and therefore closed, and it contains
$R(U)=U$; it consequently contains $K_0=\cl U$. Thus $R(K_0)=K_0$. Since
$R=AB=BA$,
\[
K_0=AB(K_0)\subseteq A(K_0)\subseteq K_0,
\qquad
K_0=BA(K_0)\subseteq B(K_0)\subseteq K_0,
\]
so both inclusions are equalities. Equation \eqref{eq:irreducible} gives $K_0=\I$. Hence
$\bigcup I_n$ is dense.

Any two continuous maps commuting with $R$ and agreeing on $I_0$ agree on every $I_n$ and hence on
$\I$. Suppose that $a^m=b^n=q$ for positive integers $m,n$. Choose a germ radius
$\varepsilon_0>0$ such that
\[
A(p+t)=p+at,\qquad B(p+t)=p+bt,\qquad R(p+t)=p+st
\quad(0\leq t\leq\varepsilon_0).
\]
After choosing
\[
0<\delta\leq
\min\{\varepsilon_0/s,\varepsilon_0/a^m,\varepsilon_0/b^n\}
\]
all intermediate iterates remain in the corresponding germ intervals, and induction gives
\[
A^m(p+t)=p+qt=B^n(p+t)\qquad(0\leq t\leq\delta).
\]
The original $A$-, $B$-, and $R$-germ estimates also hold on $J_0=[p,p+\delta]$. Repeating the
returned-hull construction with $J_0$ in place of $I_0$ shows that
$\bigcup_{k\geq0}R^k(J_0)$ is dense. The maps $A^m$ and $B^n$ commute with $R$ and agree on $J_0$,
so they agree on this dense union and hence, by continuity, on all of $\I$. Thus
$P^{Lm}=Q^{Ln}$, contradicting Lemma~\ref{lem:common-iterate}. Therefore $a,b$ are multiplicatively
independent.

For each $t\geq0$, choose any $n$ such that $t/s^n\leq\varepsilon$ and define
\begin{equation}\label{eq:Psi}
\Psi(t)=R^n\left(p+\frac{t}{s^n}\right).
\end{equation}
This is independent of the admissible choice. Indeed, $R(p+x)=p+sx$ in the chosen germ, so
consecutive admissible representatives agree:
\[
R^{n+1}\left(p+\frac{t}{s^{n+1}}\right)
=R^n\left(R\left(p+\frac{t}{s^{n+1}}\right)\right)
=R^n\left(p+\frac{t}{s^n}\right).
\]
Iteration gives agreement for any two admissible choices. In particular, for $0\leq t\leq\varepsilon$
we may take $n=0$, and hence $\Psi(t)=p+t$.
On each compact parameter interval
the same representative works throughout: if $[u,v]\subset[0,\infty)$ is compact, choose $n$ with
$v<s^n\varepsilon$. Then
\[
t\longmapsto R^n\left(p+\frac{t}{s^n}\right)
\]
defines $\Psi$ on all of $[u,v]$. Pull a finite affine partition for $R^n$ back by
$t=s^n(x-p)$. Every affine branch of $R^n$ has slope of absolute value $s^n$, so every pulled-back
branch of $\Psi$ has slope of absolute value one. This proves, uniformly on compact parameter intervals,
continuity and locally finite-piece affine unit speed. Equations \eqref{eq:germs} and commutation give \eqref{eq:equivariance}, while
$\Psi([0,s^n\varepsilon])=I_n$, so the image is dense.

It remains to prove eventual periodicity. By Lemma~\ref{lem:finite-postcritical}, the set
$\mathcal C$ in \eqref{eq:postcritical} is finite.

Put $\mathcal K=\mathcal C\cup\{p\}$. It is a finite forward-invariant set containing all turns and
endpoints of $A,B$. On a component of $\I\setminus\mathcal K$ neither map changes monotonicity; since
its adjacent affine slopes have the same prescribed magnitude $a$ (respectively $b$), continuity also
rules out a same-direction kink. Thus both maps are affine on every complementary edge, and
$\mathcal K$ is a common Markov partition. Give both orientations of each Markov edge as an alphabet.
The endpoints of an edge map into $\mathcal K$, so the oriented image under $A$ or $B$, cut at the
intervening Markov vertices, is a nonempty word of directed Markov edges. This defines nonerasing
substitutions $\Sigma_A,\Sigma_B$.
If $|e|$ is edge length, then
\begin{equation}\label{eq:length-sub}
\sum_{f\in\Sigma_A(e)}|f|=a|e|,
\qquad
\sum_{f\in\Sigma_B(e)}|f|=b|e|.
\end{equation}

The fold $\Psi$ has a literal edge itinerary $x$. Iterating \eqref{eq:vertices} gives
\[
\mathcal V(R^N)=\bigcup_{j=0}^{N-1}R^j\mathcal V(R)
 =\bigcup_{j=1}^{N}R^j(E(R)).
\]
Since $\mathcal V(R)=\mathcal V(A)\cup A(\mathcal V(B))\subseteq\mathcal C$, every value of
$R^N$ at a turn of $R^N$ belongs to $\mathcal K$. Because $R^N$ has constant nonzero slope magnitude
$s^N$ on its affine pieces, every genuine affine breakpoint is such a turn. Formula \eqref{eq:Psi}
therefore shows that every
fold value of $\Psi$ is a Markov vertex. Cut the parameter path both at these folds and whenever an
affine piece of $\Psi$ crosses a Markov vertex. The resulting pieces are successive traversals of directed
Markov edges. They cannot accumulate in bounded parameter time: if $\delta>0$ is the minimum length
of a Markov edge, every completed traversal takes exactly its edge length, hence at least $\delta$, because
$\Psi$ has unit speed.

Let $k_+$ be the least point of $\mathcal K$ greater than $p$, let $e_+$ be the directed edge from
$p$ to $k_+$, and put $\ell=k_+-p$. Since $\Psi(t)=p+t$ initially and every fold value is a Markov vertex, there can be neither a fold nor a
Markov-vertex crossing before the first point $k_+>p$ of $\mathcal K$. Hence the itinerary of $\Psi$
on $[0,\ell]$ is exactly $e_+$. The positive germs \eqref{eq:germs} make both
$\Sigma_A(e_+)$ and $\Sigma_B(e_+)$ begin with $e_+$.

For any oriented Markov edge $e$, the restriction of $A$ to its underlying interval is affine and
strictly monotone. Consequently the traversal of $A(e)$ is exactly the directed-edge concatenation
$\Sigma_A(e)$, with no backtracking; the same holds for $B$. Equation \eqref{eq:length-sub} aligns
the block-boundary times of these traversals with the parameter dilations $t\mapsto at$ and
$t\mapsto bt$. Concatenation gives the corresponding assertion for every finite word. A literal
unit-speed directed-edge itinerary is unique away from its discrete cut times, and continuity identifies
the adjacent pieces at a cut. Since
$\Psi(a^nt)=A^n\Psi(t)$ and $\Psi(b^nt)=B^n\Psi(t)$, induction therefore gives, for every $n\ge0$,
\[
\Sigma_A^n(e_+)=\text{the itinerary of $\Psi$ on }[0,a^n\ell],\qquad
\Sigma_B^n(e_+)=\text{the itinerary of $\Psi$ on }[0,b^n\ell].
\]
As $a,b>1$, these prefixes exhaust the infinite itinerary, so
\[
x=\Sigma_A^\omega(e_+)=\Sigma_B^\omega(e_+).
\]
Let $\mathcal A$ be the set of directed edges that occur in $x$. It is invariant under both substitutions:
if $e\in\mathcal A$, then every letter of $\Sigma_A(e)$ occurs in
$\Sigma_A(x)=x$, and similarly for $B$. Hence both substitutions restrict to nonerasing substitutions on
$\mathcal A$, still prolongable on $e_+$, and every letter of $\mathcal A$ occurs in their common fixed
word $x$. The restrictions are growing: if $L_{\max}$ is the largest Markov-edge length, then
\eqref{eq:length-sub} gives
$\#\Sigma_A^n(e)\ge a^n|e|/L_{\max}\to\infty$, and similarly for $B$. With the incidence convention
$M_{e,f}=\#\{f\text{ in }\Sigma(e)\}$ and with $\boldsymbol\ell>0$ the vector of lengths of the
retained edges, \eqref{eq:length-sub} gives
\[
M_A\boldsymbol\ell=a\boldsymbol\ell,\qquad
M_B\boldsymbol\ell=b\boldsymbol\ell.
\]
Durand uses the transpose, column-source incidence convention; transposition does not change spectral
radius. In either convention, Collatz--Wielandt shows that the spectral radii, hence the dominant
eigenvalues in Durand's terminology, are $a,b$.

Durand's growing-substitution theorem \cite[Theorem~9]{Durand} now applies: two nonerasing growing
substitutions, prolongable on the same letter, with every alphabet letter occurring in both common fixed
words and with multiplicatively independent dominant eigenvalues, have an ultimately periodic common
fixed word. Hence $x$ is ultimately periodic. Write $x=x_0x_1\cdots$ and put
\[
t_n=\sum_{j<n}|x_j|.
\]
If $x_{n+p}=x_n$ for every $n\geq N$, then
\[
\tau=\sum_{j=N}^{N+p-1}|x_j|>0
\quad\text{and}\quad
t_{n+p}=t_n+\tau\qquad(n\geq N).
\]
For $t\in[t_n,t_{n+1}]$, the value $\Psi(t)$ is determined by the directed edge $x_n$ and the
within-edge distance $t-t_n$. The same directed edge occurs on
$[t_n+\tau,t_{n+1}+\tau]$ with the same within-edge distance. Thus
$\Psi(t+\tau)=\Psi(t)$ for every $t\geq t_N$, so $\Psi$ is genuinely eventually periodic.
\end{proof}

\section{From a periodic fold to a fixed point}

\begin{lemma}[Periodic-path rigidity]\label{lem:path-rigidity}
Let $u,v:\mathbb R\to\I$ be periodic locally unit-speed polygonal paths satisfying
\[
u(rt)=T(u(t)),\qquad v(rt)=T(v(t))\qquad(t\in\mathbb R)
\]
for a continuous interval map $T$ and a number $r>1$.
If $v(t)=u(\varepsilon t+\delta)$ on a nonempty open interval for some
$\varepsilon\in\{\pm1\}$, then there is $\delta_0\in\mathbb R$ such that
\[
v(t)=u(\varepsilon t+\delta_0)\qquad(t\in\mathbb R).
\]
If $u$ and $v$ are onto, such a local relation always exists; consequently the two paths differ globally
by a translation or reflection, and their full period groups coincide.
\end{lemma}

\begin{proof}
Suppose first that $v(t)=u(\varepsilon t+\delta)$ on a nonempty open interval. Restricting if necessary,
fix a bounded nonempty open subinterval $I$ on which the identity holds. Equivariance gives
\begin{equation}\label{eq:expanded-local}
v(t)=u(\varepsilon t+r^n\delta)\qquad(t\in r^nI).
\end{equation}
Let $m_n$ be the midpoint of $r^nI$, and choose positive periods $p_v,p_u$ of $v,u$. Choose integers
$k_n,l_n$ and representatives
\[
a_n=m_n-k_np_v\in[0,p_v),\qquad
b_n=\varepsilon m_n+r^n\delta-l_np_u\in[0,p_u).
\]
Passing to a subsequence, $a_n\to a_0$ and $b_n\to b_0$. Since $|r^nI|\to\infty$, for each
fixed $h\in\mathbb R$ the point $m_n+h$ lies in $r^nI$ for all sufficiently large $n$. Applying
\eqref{eq:expanded-local} there and using periodicity gives
\[
v(a_n+h)=u(b_n+\varepsilon h),
\]
where $a_n\to a_0$ and $b_n\to b_0$. Continuity yields
$v(a_0+h)=u(b_0+\varepsilon h)$ for every $h\in\mathbb R$, which is the desired global relation after
putting $\delta_0=b_0-\varepsilon a_0$.

Now suppose that $u,v$ are onto. In one period each polygonal path has only finitely many vertex values.
Choose $y\in(0,1)$ outside the union of those two finite sets and choose preimages $t_u,t_v$ with
$u(t_u)=v(t_v)=y$. Neither preimage is a vertex. On sufficiently small neighbourhoods the two paths are
affine with slopes in $\{\pm1\}$, so after a translation and, if their slopes have opposite signs, a
reflection of the parameter, one obtains a local relation of the required form. The first part globalizes it.
Finally, if $v(t)=u(\varepsilon t+\delta_0)$ globally, then
$\tau\in\operatorname{Per}(v)$ exactly when $\varepsilon\tau\in\operatorname{Per}(u)$; a period group is
closed under negation, so the two full period groups are equal.
\end{proof}

\begin{lemma}[Periodic fold]\label{lem:periodic-fold}
Let $P,Q:\I\to\I$ be commuting onto flat-free finite-piece PL maps with constant absolute slopes
$C,D>1$. Let $L\geq1$, put $A=P^L$, $B=Q^L$, $a=C^L$, and $b=D^L$, and assume the returned-power irreducibility
condition \eqref{eq:irreducible}. If there is a nonconstant continuous eventually periodic locally
finite-piece affine unit-speed map $\Psi:[0,\infty)\to\I$ satisfying
\[
\Psi(at)=A\Psi(t),\qquad \Psi(bt)=B\Psi(t),
\]
then $P,Q$ have a common fixed point.
\end{lemma}

\begin{proof}
Let $\Psi(t+\tau_*)=\Psi(t)$ for $t\ge T_0$, where $\tau_*>0$. For $t\in\mathbb R$, choose an
integer $m\ge0$ with $t+m\tau_*\ge T_0$ and define
\[
\psi(t)=\Psi(t+m\tau_*).
\]
If $m,m'$ are admissible and $m'\geq m$, repeated use of eventual $\tau_*$-periodicity gives
$\Psi(t+m'\tau_*)=\Psi(t+m\tau_*)$, so the definition is independent of the choice. Given
$t_0\in\mathbb R$, choose one $m$ and a neighbourhood $W$ of $t_0$ such that
$W+m\tau_*\subset[T_0,\infty)$. On $W$ the map $\psi$ is the translate
$t\mapsto\Psi(t+m\tau_*)$ and hence is continuous and locally finite-piece affine of unit speed. A compact
interval has a finite cover by such neighbourhoods; refining at the finitely many cover and affine-piece
endpoints gives a finite unit-speed partition. Unit speed on any nondegenerate compact interval in the
tail also shows that the tail, and hence $\psi$, is nonconstant. Thus $\psi$ is a continuous periodic locally
unit-speed polygonal path extending the periodic tail. Its full period group
\[
\operatorname{Per}(\psi)=\{\tau\in\mathbb R:\psi(t+\tau)=\psi(t)\ \text{for all }t\}
\]
is a closed subgroup of $\mathbb R$. It contains $\tau_*\ne0$, and it is not all of $\mathbb R$ because
$\psi$ is nonconstant; hence the classification of closed subgroups of $\mathbb R$ gives
$\operatorname{Per}(\psi)=\tau_0\mathbb Z$ for a unique $\tau_0>0$. For large $t$,
\[
\psi(at)=A\psi(t),\qquad \psi(bt)=B\psi(t).
\]
For all sufficiently large $t$,
\[
\psi(a(t+\tau_0))=A\psi(t+\tau_0)=A\psi(t)=\psi(at),
\]
so $a\tau_0$ is an eventual period. Given any $u\in\mathbb R$, translate both $u$ and $u+a\tau_0$
by a large integer multiple of the already global period $\tau_0$; the eventual identity then shows that
$a\tau_0$ is a global period. Thus $a\tau_0\in\tau_0\mathbb Z$, so $a\in\mathbb N$. Similarly
$b\in\mathbb N$. Now, for arbitrary $t$, choose an integer $M$ with $t+M\tau_0$ in the range where
$\psi(at)=A\psi(t)$ is known. Since $aM\tau_0$ is a period,
\[
\psi(at)=\psi(a(t+M\tau_0))=A\psi(t+M\tau_0)=A\psi(t),
\]
and similarly for $B$. Thus both return identities are global. The image
$J=\psi(\mathbb R)=\psi([0,\tau_0])$ is compact by periodicity, connected because $\mathbb R$ is
connected, and nondegenerate because $\psi$ is nonconstant. The global identities give
$A(J)=J=B(J)$, so \eqref{eq:irreducible} gives $J=\I$; hence $\psi$ is onto.

Define
\[
v_P(t)=P\psi(t/C),\qquad v_Q(t)=Q\psi(t/D).
\]
They are onto because $P,Q$, and $\psi$ are onto, and $C\tau_0,D\tau_0$ are respective periods.
To check local unit speed for $v_P$, start with a compact parameter
interval, take a finite unit-speed affine partition for $\psi$ on the corresponding interval scaled by
$1/C$, and pull it back by $t\mapsto t/C$. Refine the pulled-back partition at the preimages of the
finitely many affine breakpoints of $P$. On every resulting cell the inner path has slope $\pm1/C$ and
the relevant branch of $P$ has slope of absolute value $C$, so their composite has slope $\pm1$.
This refinement is finite. The argument for $v_Q$ is identical, with $D$ in place of $C$. Using the
global return identities and commutation,
\[
v_P(at)=P\psi(at/C)=PA\psi(t/C)=A v_P(t),
\qquad
v_Q(bt)=B v_Q(t).
\]
Lemma~\ref{lem:path-rigidity} gives
\[
v_P(t)=\psi(\varepsilon_Pt+\delta_P),\qquad
v_Q(t)=\psi(\varepsilon_Qt+\delta_Q).
\]
Since the full period groups coincide, the periods $C\tau_0$ of $v_P$ and $D\tau_0$ of $v_Q$ belong
to $\tau_0\mathbb Z$. Hence $C,D\in\mathbb N$. Substituting $t=Cx,Dx$ produces affine maps of the
parameter circle $\mathbb R/\tau_0\mathbb Z$,
\begin{equation}\label{eq:lifts}
P\psi(x)=\psi(kx+\alpha),\qquad Q\psi(x)=\psi(\ell x+\beta),
\end{equation}
with nonzero integers $k=\varepsilon_PC$, $\ell=\varepsilon_QD$. These affine circle maps commute.
Indeed, $PQ=QP$ and \eqref{eq:lifts} give
\[
\psi(k\ell t+k\beta+\alpha)=\psi(k\ell t+\ell\alpha+\beta)\qquad(t\in\mathbb R).
\]
Put $\Delta=(k\beta+\alpha)-(\ell\alpha+\beta)$. On the parameter circle
$\mathbb R/\tau_0\mathbb Z$, multiplication by the nonzero integer $k\ell$ is surjective. Therefore, for
every $y$ on that circle, one can choose $t$ with
$y=k\ell t+\ell\alpha+\beta$, and the preceding identity becomes
$\psi(y+\Delta)=\psi(y)$. Hence $\Delta\in\operatorname{Per}(\psi)=\tau_0\mathbb Z$, which is
precisely $\mathcal P\mathcal Q=\mathcal Q\mathcal P$ on the parameter circle.

Choose $t_0$ with $\psi(t_0)=1$. Since $\psi\le1$, the point $t_0$ cannot lie in the interior of an
affine unit-slope piece. The two one-sided slopes at this maximum are therefore $+1$ and $-1$, so
$\psi(t_0-h)=1-h=\psi(t_0+h)$ for all sufficiently small $h\ge0$. Applying
Lemma~\ref{lem:path-rigidity} to $u=v=\psi$ globalizes this local symmetry to a reflection
\[
\rho(t)=-t+\gamma,\qquad \psi\rho=\psi.
\]
It is unique modulo periods, because the product of two reflection symmetries is a translation symmetry.
The affine lifts commute with $\rho$. For example, write
$\mathcal P(t)=kt+\alpha$ on $\mathbb R/\tau_0\mathbb Z$. From
\[
\psi(\mathcal P\rho(t))=P\psi(\rho(t))=P\psi(t)=\psi(\mathcal P(t))
\]
and the surjectivity of $\mathcal P$, putting $y=\mathcal P(t)$ gives
\[
\psi(-y+k\gamma+2\alpha)=\psi(y)\qquad(y\in\mathbb R/\tau_0\mathbb Z).
\]
Thus $y\mapsto-y+k\gamma+2\alpha$ is a reflection symmetry of $\psi$. By uniqueness it equals
$\rho$, which is equivalent to $\mathcal P\rho=\rho\mathcal P$. The same argument applies to
$\mathcal Q$.

Rescale the parameter circle to $\mathbb R/\mathbb Z$ and translate the origin to a fixed point of
$\rho$; retaining the same notation for the reparametrized objects, we have $\rho(t)=-t$. Write the
affine lifts as
\[
\mathcal P(t)=kt+\alpha,\qquad \mathcal Q(t)=\ell t+\beta.
\]
Commutation with $\rho$ gives
$\alpha,\beta\in\{0,\tfrac12\}$, and affine commutation gives
\begin{equation}\label{eq:affine-comm}
(k-1)\beta=(\ell-1)\alpha\pmod1.
\end{equation}
On the four-torsion
$\{0,\tfrac14,\tfrac12,\tfrac34\}$ consider the three reflection orbits
\[
\{0\},\qquad\{\tfrac14,\tfrac34\},\qquad\{\tfrac12\}.
\]
There is always a common fixed orbit. If $\alpha=\beta=0$, use $\{0\}$. If
$\alpha=0,\beta=\tfrac12$, \eqref{eq:affine-comm} makes $k$ odd, so $\mathcal P$ fixes all three orbits;
$\mathcal Q$ fixes the middle orbit when $\ell$ is odd and $\{\tfrac12\}$ when $\ell$ is even. The
case $\alpha=\tfrac12,\beta=0$ is symmetric. If $\alpha=\beta=\tfrac12$, then $k,\ell$ have the same
parity; both fix the middle orbit when odd and $\{\tfrac12\}$ when even.

Choose $t$ in a common fixed reflection orbit. Then
$\mathcal P(t),\mathcal Q(t)\in\{t,-t\}$ and $\psi(-t)=\psi(t)$. Equation \eqref{eq:lifts} gives
$P\psi(t)=\psi(t)=Q\psi(t)$.
\end{proof}

\begin{proof}[Proof of Theorem~\ref{thm:main}]
Assume a counterexample exists. Lemma~\ref{lem:linearize} gives a minimal expanding constant-slope
counterexample. Lemma~\ref{lem:irreducible} gives an extremal factor satisfying
\eqref{eq:irreducible}. Lemma~\ref{lem:two-scale}, after at most a simultaneous reflection conjugacy,
gives an eventually periodic two-scale fold. Lemma~\ref{lem:periodic-fold} gives a common fixed point of
the conjugated extremal pair. Conjugating that point back if a reflection was used gives a common fixed
point of the original extremal factor, a contradiction.
\end{proof}

\begin{thebibliography}{99}
\bibitem{Boyce} W.~M. Boyce, \emph{Commuting functions with no common fixed point},
Trans. Amer. Math. Soc. \textbf{137} (1969), 77--92.
\bibitem{Huneke} J.~P. Huneke, \emph{On common fixed points of commuting continuous functions on an interval},
Trans. Amer. Math. Soc. \textbf{139} (1969), 371--381.
\bibitem{Durand} F.~Durand, \emph{Cobham's theorem for substitutions},
J. Eur. Math. Soc. \textbf{13} (2011), 1799--1814.
\bibitem{AlsedaMisiurewicz} Ll.~Alsed\`a and M.~Misiurewicz,
\emph{Semiconjugacy to a map of a constant slope}, Discrete Contin. Dyn. Syst. Ser. B \textbf{20}
(2015), 3403--3413.
\bibitem{FugateMcLean} J.~B. Fugate and T.~B. McLean,
\emph{Compact groups of homeomorphisms on tree-like continua},
Trans. Amer. Math. Soc. \textbf{267} (1981), 609--620,
doi: \href{https://doi.org/10.1090/S0002-9947-1981-0626493-5}
{10.1090/S0002-9947-1981-0626493-5}.
\bibitem{BogatyiFrolkina} S.~A. Bogatyi and O.~D. Frolkina,
\emph{A common fixed point of commuting mappings of a tree}, Vestnik Moskov. Univ. Ser.~1 Mat. Mekh.
2002, no.~6, 3--10; English transl., Moscow Univ. Math. Bull. \textbf{57} (2002), no.~6, 1--8 (2003).
\end{thebibliography}

\end{document}
